\documentclass{article}
\usepackage[T1]{fontenc}
\usepackage{lmodern}
\usepackage{graphicx}
\usepackage{amsmath,amssymb}
\usepackage[a4paper,margin=2cm]{geometry}
\usepackage[absolute,overlay]{textpos}
\setlength{\TPHorizModule}{1mm}
\setlength{\TPVertModule}{1mm}
\usepackage[indonesian]{babel}
\usepackage{hyperref}
\setlength{\emergencystretch}{2em}
% unit: Halaman 48

\begin{document}

% === Halaman 48 (python/native) ===

48

Metode Kasiski

• Metode Kasiski tidak secara langsung menemukan kunci Vigenere Cipher, 

tetapi membantu menemukan estimasi panjang kunci Vigenere cipher.

• Metode Kasiski memanfaatkan keuntungan bahwa bahasa Inggris tidak hanya 

mengandung perulangan huruf,

• tetapi juga perulangan pasangan huruf atau tripel huruf, seperti TH, THE, EN,  

dsb.

• Perulangan kelompok huruf ini ada kemungkinan menghasilkan kriptogram 

yang berulang.

\end{document}
